%----------------------------------------------------------------------------------------
%       ANEXO A: MANUAL DE USUARIO
%----------------------------------------------------------------------------------------
\chapter{Manual de Usuario - Interfaz Gráfica de Demostración}
\label{anexo:manual_usuario}

Este anexo presenta el manual de usuario de la interfaz gráfica desarrollada para demostrar el sistema completo de detección de COVID-19 mediante landmarks anatómicos y normalización geométrica.

\section{Descripción General}

La interfaz gráfica está implementada con Gradio, un framework para crear interfaces web interactivas en Python. El sistema permite visualizar las cuatro etapas principales del pipeline de procesamiento:

\begin{enumerate}
    \item Imagen original de radiografía de tórax
    \item Detección de 15 landmarks anatómicos
    \item Imagen normalizada mediante warping geométrico
    \item Mapa de atención GradCAM mostrando regiones relevantes para la clasificación
\end{enumerate}

\section{Requisitos del Sistema}

\subsection{Dependencias de Software}

\begin{itemize}
    \item Python 3.8 o superior
    \item Gradio 4.0.0 o superior
    \item PyTorch con soporte para GPU (opcional, recomendado)
    \item Dependencias adicionales especificadas en \texttt{requirements.txt}
\end{itemize}

\subsection{Modelos Necesarios}

El sistema requiere los siguientes archivos de modelos entrenados:

\textbf{Ensemble de Landmarks (4 modelos):}
\begin{itemize}
    \item \texttt{checkpoints/session10/ensemble/seed123/final\_model.pt}
    \item \texttt{checkpoints/session13/seed321/final\_model.pt}
    \item \texttt{checkpoints/repro\_split111/session14/seed111/final\_model.pt}
    \item \texttt{checkpoints/repro\_split666/session16/seed666/final\_model.pt}
\end{itemize}

\textbf{Forma Canónica y Triangulación:}
\begin{itemize}
    \item \texttt{outputs/shape\_analysis/canonical\_shape\_gpa.json}
    \item \texttt{outputs/shape\_analysis/canonical\_delaunay\_triangles.json}
\end{itemize}

\textbf{Clasificador:}
\begin{itemize}
    \item {\small \texttt{outputs/classifier\_warped\_lung\_best/sweeps\_2026-01-12/\linebreak lr2e-4\_seed321\_on/best\_classifier.pt}}
\end{itemize}

\section{Instalación y Ejecución}

\subsection{Lanzar la Interfaz}

Existen dos opciones para ejecutar la interfaz:

\paragraph{Opción 1 (Recomendada):} Usar el script de lanzamiento con verificaciones automáticas:

\begin{verbatim}
python scripts/run_demo.py
\end{verbatim}

Opciones disponibles:
\begin{itemize}
    \item \texttt{-{}-share}: Crear enlace público compartible
    \item \texttt{-{}-port PORT}: Cambiar puerto (predeterminado: 7860)
    \item \texttt{-{}-host HOST}: Cambiar host (predeterminado: localhost)
\end{itemize}

\paragraph{Opción 2:} Ejecutar directamente el módulo:

\begin{verbatim}
python -m src_v2.gui.app
\end{verbatim}

La interfaz se abrirá automáticamente en el navegador web en \texttt{http://localhost:7860}.

\section{Uso de la Interfaz}

\subsection{Tab 1: Demostración Completa}

Esta pestaña muestra el pipeline completo con visualizaciones de todas las etapas:

\begin{enumerate}
    \item \textbf{Cargar imagen:}
    \begin{itemize}
        \item Arrastrar y soltar una radiografía de tórax en el área designada
        \item Hacer clic en el botón de carga para seleccionar un archivo
        \item Seleccionar uno de los ejemplos precargados (COVID, Normal, Viral)
    \end{itemize}

    \item \textbf{Procesar:}
    \begin{itemize}
        \item Hacer clic en el botón ``Procesar Imagen''
        \item Esperar 1-2 segundos mientras el sistema procesa (el tiempo varía según el hardware)
    \end{itemize}

    \item \textbf{Visualizar resultados:}
    \begin{itemize}
        \item La interfaz mostrará 4 imágenes correspondientes a cada etapa del pipeline
        \item Se mostrarán las probabilidades de clasificación para cada clase
        \item Es posible expandir la sección ``Métricas Detalladas'' para ver el error promedio por landmark
    \end{itemize}

    \item \textbf{Exportar (opcional):}
    \begin{itemize}
        \item Hacer clic en ``Exportar Resultados a PDF''
        \item El sistema generará un reporte PDF con todas las visualizaciones y métricas
        \item El archivo se guardará en el directorio de trabajo actual
    \end{itemize}
\end{enumerate}

\subsection{Tab 2: Vista Rápida}

Esta pestaña proporciona una clasificación directa sin visualizaciones intermedias:

\begin{enumerate}
    \item Cargar la radiografía de tórax
    \item Hacer clic en ``Clasificar''
    \item Obtener el resultado inmediato (clase predicha y probabilidades)
\end{enumerate}

Esta opción es útil cuando se requiere únicamente la clasificación sin inspeccionar las etapas intermedias.

\subsection{Tab 3: Acerca del Sistema}

Esta pestaña contiene información sobre:
\begin{itemize}
    \item Metodología de landmarks y normalización geométrica
    \item Arquitectura del sistema
    \item Métricas validadas del modelo
    \item Referencias bibliográficas
\end{itemize}

\section{Visualización de Landmarks}

Los 15 landmarks se visualizan con un código de colores según su grupo anatómico:

\begin{table}[H]
\centering
\begin{tabular}{llp{6cm}}
\toprule
\textbf{Grupo} & \textbf{Landmarks} & \textbf{Descripción} \\
\midrule
Eje & L1, L2 & Puntos superior e inferior del eje central \\
Central & L9, L10, L11 & Puntos intermedios del eje central \\
Lateral & L3--L8 & Contornos laterales izquierdo y derecho \\
Borde & L12, L13 & Puntos de borde superior \\
Costal & L14, L15 & Puntos costales inferiores \\
\bottomrule
\end{tabular}
\caption{Grupos de landmarks y su representación visual}
\end{table}

\section{Métricas del Sistema}

Las métricas validadas que se muestran en la interfaz son:

\begin{table}[H]
\centering
\begin{tabular}{lc}
\toprule
\textbf{Métrica} & \textbf{Valor} \\
\midrule
Error de Landmarks (píxeles) & $3.61 \pm 2.48$ \\
Accuracy de Clasificación & 98.05\% \\
F1-Score Macro & 97.12\% \\
F1-Score Weighted & 98.04\% \\
\bottomrule
\end{tabular}
\caption{Métricas validadas del sistema completo}
\end{table}

Estos valores fueron obtenidos mediante evaluación rigurosa en el conjunto de prueba y están documentados en el archivo \texttt{GROUND\_TRUTH.json} versión 2.1.0.

\section{Solución de Problemas}

\subsection{Error: Modelos no encontrados}

Si la interfaz reporta que no encuentra los modelos:
\begin{itemize}
    \item Verificar que los archivos de checkpoints existen en las rutas especificadas
    \item Ejecutar \texttt{python scripts/run\_demo.py} para diagnóstico automático
    \item Consultar la documentación de reproducción del pipeline completo
\end{itemize}

\subsection{Error: Memoria GPU insuficiente}

La interfaz automáticamente detectará problemas de memoria GPU y realizará fallback a CPU:
\begin{itemize}
    \item El tiempo de inferencia será aproximadamente 2-3× más lento en CPU
    \item Para forzar el uso de CPU, modificar \texttt{DEVICE\_PREFERENCE = "cpu"} en \texttt{config.py}
\end{itemize}

\subsection{La interfaz no se abre en el navegador}

\begin{itemize}
    \item Verificar que el puerto 7860 no esté siendo utilizado por otra aplicación
    \item Usar la opción \texttt{-{}-port 8080} para cambiar el puerto
    \item Revisar la configuración del firewall si se usa la opción \texttt{-{}-share}
\end{itemize}

\subsection{Imágenes de baja calidad}

\begin{itemize}
    \item Tamaño mínimo recomendado: 100×100 píxeles
    \item Tamaño óptimo: 224×224 píxeles o superior
    \item La interfaz redimensiona automáticamente las imágenes al tamaño requerido
\end{itemize}

\section{Arquitectura del Sistema}

\subsection{Patrón Singleton para Gestión de Modelos}

El módulo \texttt{ModelManager} utiliza el patrón Singleton para optimizar el uso de memoria:
\begin{itemize}
    \item \textbf{Lazy loading}: Los modelos se cargan únicamente cuando se utilizan por primera vez
    \item \textbf{Cacheo en memoria}: Los modelos permanecen en memoria entre inferencias
    \item \textbf{Detección automática}: El sistema detecta automáticamente si hay GPU disponible
\end{itemize}

\subsection{Pipeline de Inferencia}

El flujo de procesamiento sigue las siguientes etapas:

\begin{enumerate}
    \item \textbf{Validación}: Verificar formato y dimensiones de la imagen
    \item \textbf{Preprocesamiento}: Cargar y redimensionar a 224×224 píxeles
    \item \textbf{Predicción de landmarks}:
    \begin{itemize}
        \item Aplicar CLAHE (clip=2.0, tile\_size=4)
        \item Ejecutar ensemble de 4 modelos
        \item Aplicar TTA (Test-Time Augmentation) con flip horizontal
        \item Promediar predicciones para obtener landmarks finales (15, 2)
    \end{itemize}
    \item \textbf{Warping}: Normalización geométrica mediante transformación afín por piezas
    \item \textbf{Clasificación}: Inferencia con ResNet-18 y generación de GradCAM
    \item \textbf{Visualización}: Renderizado de las 4 etapas para presentación al usuario
\end{enumerate}

\section{Notas Técnicas}

\subsection{Formato de Imágenes Soportadas}

La interfaz acepta los siguientes formatos:
\begin{itemize}
    \item PNG (.png)
    \item JPEG (.jpg, .jpeg)
    \item Escala de grises o RGB (se convierte automáticamente a escala de grises)
\end{itemize}

\subsection{Tiempo de Procesamiento}

El tiempo de procesamiento típico es:
\begin{itemize}
    \item Con GPU (NVIDIA con CUDA): 1-2 segundos
    \item Con CPU (multi-core moderno): 3-5 segundos
    \item El primer procesamiento puede tardar más debido a la carga inicial de modelos
\end{itemize}

\subsection{Uso de Recursos}

\begin{itemize}
    \item Memoria GPU requerida: $\sim$2 GB
    \item Memoria RAM requerida: $\sim$4 GB
    \item Espacio en disco para modelos: $\sim$500 MB
\end{itemize}
